% Part: computability
% Chapter: recursive-functions
% Section: sequences

\documentclass[../../../include/open-logic-section]{subfiles}

\begin{document}

\olfileid{cmp}{rec}{seq}
\olsection{અનુક્રમો}

આદિમ પુનરાવર્તી વિધેયોનો ગણ ખાસ્સો સમર્થ છે. પરંતુ
\emph{અનુક્રમો}ને સંભાળવાની યોગ્ય રીત વિકસાવીએ પછી આપણે વધુ કરી
શકીશું. પ્રાકૃતિક સંખ્યાઓના સાન્ત અનુક્રમોને પ્રાકૃતિક સંખ્યાઓ સાથે
આ રીતે ઓળખીશું: અનુક્રમ $\langle a_0, a_1, a_2, \dots, a_k \rangle$
નીચેની સંખ્યાને અનુરૂપ છે:
\[
p_0^{a_0+1} \cdot p_1^{a_1+1} \cdot p_2^{a_2+1} \cdot \dots \cdot
p_k^{a_k+1}.
\]
ઘાતમાં એક ઉમેરીએ છીએ, જેથી દાખલા તરીકે
$\langle 2, 7, 3\rangle$ અને $\langle 2, 7, 3, 0, 0 \rangle$ના
સંખ્યાત્મક કોડ જુદા રહે. ખાલી અનુક્રમના કોડ તરીકે $0$ અને~$1$ બંને
લઈ શકાય; ચોક્કસતા માટે $\emptyseq$ એ~$0$ દર્શાવે એમ માનીએ.

અનુક્રમોનું આ સંકેતીકરણ સફળ થાય છે તેનું કારણ અંકગણિતનું કહેવાતું
મૂળભૂત પ્રમેય છે: દરેક પ્રાકૃતિક સંખ્યા $n \ge 2$ને નીચેના
સ્વરૂપમાં એક જ રીતે લખી શકાય:
\[
n = p_0^{a_0} \cdot p_1^{a_1} \cdot \dots \cdot p_k^{a_k}
\]
અહીં $a_k \ge 1$ છે. આથી નિશ્ચિત થાય છે કે વિધેય
$\tuple{}(a_0,
\dots, a_k) = \tuple{a_0, \dots, a_k}$ એક-એક છે: જુદા
અનુક્રમો જુદી સંખ્યાઓ પર જાય છે; દરેક સંખ્યાને વધુમાં વધુ એક
અનુક્રમ અનુરૂપ છે.

હવે બતાવીશું કે અનુક્રમની લંબાઈ શોધવી, તેનો $i$મો ઘટક શોધવો,
તેમાં ઘટક અંતે ઉમેરવો અને બે અનુક્રમો જોડવા---આ બધી ક્રિયાઓ આદિમ
પુનરાવર્તી છે.

\begin{prop}
  અનુક્રમ $s$ની લંબાઈ આપતું વિધેય $\len{s}$ આદિમ પુનરાવર્તી છે.
\end{prop}

\begin{proof}
  સંબંધ $R(i, s)$ને આ રીતે વ્યાખ્યાયિત કરીએ:
  \[
  R(i, s) \text{ ત્યારે અને માત્ર ત્યારે જ જ્યારે }
  p_i \mid s \land p_{i+1} \nmid s.
  \]
  સ્પષ્ટ છે કે $R$ આદિમ પુનરાવર્તી છે. જ્યારે પણ $s$ એ બિનખાલી
  અનુક્રમનો કોડ હોય, એટલે
  \[
  s = p_0^{a_0+1} \cdot \dots \cdot p_{k}^{a_{k}+1},
  \]
  ત્યારે $p_i$ એ એવી સૌથી મોટી અવિભાજ્ય સંખ્યા હોય કે
  $p_i \mid s$, તો $R(i,s)$ લાગુ પડે છે, એટલે $i = k$.
  તેથી $p_i$ એ~$s$ને ભાગતી
  સૌથી મોટી અવિભાજ્ય સંખ્યા હોય ત્યારે અને માત્ર ત્યારે જ $s$ની
  લંબાઈ $i+1$ છે. આથી આમ રાખી શકીએ:
  \[
  \len{s} =
  \begin{cases}
    0 & \text{જો $s = 0$ અથવા $s = 1$ હોય} \\
    1 + \bmin{i < s}{R(i, s)} & \text{અન્યથા}
  \end{cases}
  \]
  અહીં સીમાબદ્ધ લઘુત્તમીકરણ વાપરી શકીએ છીએ, કારણ કે $s$ એ અનુક્રમનો
  કોડ હોય ત્યારે $R(i,s)$ સંતોષતો એક જ $i$ હોય છે; અને એવો $i$ હોય તો
  પોતે~$s$થી નાનો છે.
\end{proof}

\begin{prop}
  અનુક્રમ $s$ના અંતે $a$ ઉમેરવાથી મળતું અનુક્રમ આપતું વિધેય
  $\fn{append}(s,a)$ આદિમ પુનરાવર્તી છે.
\end{prop}

\begin{proof}
  $\fn{append}$ને આમ વ્યાખ્યાયિત કરી શકીએ:
  \[
  \fn{append}(s,a) =
  \begin{cases}
    2^{a+1} & \text{જો $s = 0$ અથવા $s = 1$ હોય} \\
    s \cdot p_{\len{s}}^{a+1} & \text{અન્યથા.}
  \end{cases}
  \]
\end{proof}

\begin{prop}
  $s$નો $i$મો ઘટક આપતું વિધેય $\fn{element}(s,i)$ આદિમ
  પુનરાવર્તી છે; અહીં શરૂઆતના ઘટકને $0$મો કહીએ છીએ. જો $i$ એ
  $s$ની લંબાઈ જેટલો અથવા તેનાથી મોટો હોય, તો વિધેય $0$ આપે છે.
\end{prop}

\begin{proof}
$a$ એ~$s$નો $i$મો ઘટક હોય ત્યારે અને માત્ર ત્યારે જ
$p_i^{a+1}$ એ~$s$ને ભાગતો $p_i$નો સૌથી મોટો ઘાત છે, એટલે
$p_i^{a+1} \mid s$ પરંતુ $p_i^{a+2} \nmid s$. તેથી:
  \[
  \fn{element}(s,i) =
  \begin{cases}
    0 & \mbox{જો $i \geq \len{s}$ હોય} \\
    \bmin{a < s}{(p_i^{a+2} \nmid s)} & \text{અન્યથા.}
  \end{cases}
  \]
\end{proof}

ઉપર વ્યાખ્યાયિત વિધેયોનાં ઔપચારિક નામોને બદલે હવે વધુ ટૂંકાં
સંકેતનો લાવીએ. $\fn{element}(s,i)$ને બદલે $(s)_i$ લખીશું, અને
નીચેના સંક્ષેપ તરીકે $\tuple{s_0, \dots, s_k}$ લખીશું:
\[
\fn{append}(\fn{append}(\dots \fn{append}(\emptyseq,s_0)
\dots),s_k).
\]
નોંધો કે $s$ની લંબાઈ~$k$ હોય, તો $s$ના ઘટકો
$(s)_0$, \dots,~$(s)_{k-1}$ છે.

\begin{prop}
બે અનુક્રમો જોડતું વિધેય $\fn{concat}(s,t)$ આદિમ પુનરાવર્તી છે.
\end{prop}

\begin{proof}
  આપણે એવું વિધેય $\fn{concat}$ જોઈએ છે કે
  \[
    \fn{concat}(\tuple{a_0, \dots, a_k}, \tuple{b_0, \dots, b_l}) =
    \tuple{a_0, \dots, a_k, b_0, \dots, b_l}.
  \]
  સહાયક વિધેય $\fn{hconcat}(s,t,n)$ વાપરીશું, જે $t$ના પહેલાં
  $n$ ઘટકોને~$s$ના અંતે જોડે છે. તેને આદિમ પુનરાવર્તન વડે આમ
  વ્યાખ્યાયિત કરી શકાય:
  \begin{align*}
    \fn{hconcat}(s,t,0) & = s\\
    \fn{hconcat}(s,t,n+1) & = \fn{append}(\fn{hconcat}(s,t,n),(t)_n)
    \intertext{પછી $\fn{concat}$ને આમ વ્યાખ્યાયિત કરીએ:}
    \fn{concat}(s,t) & = \fn{hconcat}(s,t,\len{t}).
  \end{align*}
\end{proof}

$\fn{concat}(s,t)$ને બદલે આપણે $s \concat t$ લખીશું.

અનુક્રમની લંબાઈ અને તેના સૌથી મોટા ઘટકના આધારે તેનો સંખ્યાત્મક
કોડ ઉપરથી સીમિત કરી શકીએ, એ આપણને ઉપયોગી થશે. ધારો કે $s$ એ
લંબાઈ~$k$ ધરાવતો અનુક્રમ છે અને તેનો દરેક ઘટક કોઈ સંખ્યા~$x$થી
નાનો અથવા સમાન છે. તો~$s$ના વધુમાં વધુ $k$ અવિભાજ્ય અવયવો છે,
જે દરેક વધુમાં વધુ~$p_{k-1}$ છે; અને $s$ના અવિભાજ્ય અવયવીકરણમાં
દરેકનો ઘાત વધુમાં વધુ $x+1$ છે. બીજા શબ્દોમાં, જો આપણે
\[
\fn{sequenceBound}(x,k) = p_{k-1}^{k \cdot (x+1)},
\]
વ્યાખ્યાયિત કરીએ, તો ઉપર વર્ણવેલા અનુક્રમ~$s$નો સંખ્યાત્મક કોડ
વધુમાં વધુ~$\fn{sequenceBound}(x,k)$ છે.
\sourcecorrection{OLCMP-007}{સ્થિર અંગ્રેજી મૂળના આ સૂત્રમાં ખાલી
અનુક્રમ માટે $k=0$ લેતાં $p_{k-1}$ અનિર્ધારિત છે. ખાલી અનુક્રમનો કોડ
$0$ છે; આ સીમા-વિધેયને $k=0$ માટે મૂલ્ય $1$ આપીએ, જેથી ખાલી
અનુક્રમનો કોડ તેના કરતાં નાનો રહે. પ્રદર્શિત સૂત્ર $k\geq 1$ માટે છે.}

અનુક્રમો પર આવી સીમા હોવાથી સીમિત શોધ વડે નવાં વિધેયો
વ્યાખ્યાયિત કરી શકીએ. દાખલા તરીકે, $\fn{concat}$ને સીમિત શોધ
વડે વ્યાખ્યાયિત કરી શકીએ. આપણે શોધતા હોઈએ તે પદાર્થની (જોડાયેલા
અનુક્રમની સંખ્યાની) આદિમ પુનરાવર્તી \emph{વર્ણનાત્મક શરત} અને
ક્યાં સુધી શોધવું તેની સીમા લખવી પૂરતી છે. નીચેનું ચાલે છે:
\begin{align*}
  \fn{concat}(s,t) = {} & \bmin{v < \fn{sequenceBound}(s+t,\len{s} +
    \len{t})}{} \\
  & \quad(\len{v} = \len{s} + \len{t} \land {}\\
    & \qquad \bforall{i < \len{s}}{((v)_i = (s)_i) \land {} \\
      & \qquad \bforall{j < \len{t}}{((v)_{\len{s}+j} = (t)_j)})}
\end{align*}

\begin{prob}
$\fn{sconcat}(s)$ એવું આદિમ પુનરાવર્તી વિધેય છે કે જેનું
નીચેનું ગુણધર્મ છે, તે બતાવો:
\[
\fn{sconcat}(\tuple{s_0, \dots, s_k}) = s_0 \concat \dots \concat s_k.
\]
\end{prob}

\begin{prob}
$\fn{tail}(s)$ એવું આદિમ પુનરાવર્તી વિધેય છે કે જેનું
નીચેનું ગુણધર્મ છે, તે બતાવો:
\begin{align*}
  \fn{tail}(\emptyseq) & = 0 \text{ અને}\\
  \fn{tail}(\tuple{s_0, \dots, s_{k}}) & = \tuple{s_1, \dots, s_{k}}.
\end{align*}
\end{prob}

\begin{prop}
  \ollabel{prop:subseq}
  $s$ના $i$મા ઘટકથી શરૂ થતો લંબાઈ~$n$નો ઉપઅનુક્રમ આપતું વિધેય
  $\fn{subseq}(s, i, n)$ આદિમ પુનરાવર્તી છે.
\end{prop}

\begin{proof}
  અભ્યાસ માટે.
\end{proof}

\begin{prob}
  \olref[cmp][rec][seq]{prop:subseq} સાબિત કરો.
\end{prob}

\end{document}
